\documentclass[12pt,a4paper]{article}

% ------------------------------------------------------------
% Paquetes
% ------------------------------------------------------------
\usepackage[utf8]{inputenc}
\usepackage[T1]{fontenc}
\usepackage[spanish,es-tabla,es-noshorthands]{babel}
\usepackage{lmodern}
\usepackage{geometry}
\usepackage{setspace}
\usepackage{parskip}
\usepackage{microtype}
\usepackage{amsmath,amssymb,bm}
\usepackage{booktabs}
\usepackage{array}
\usepackage{siunitx}
\usepackage{xcolor}
\usepackage{float}
\usepackage[section]{placeins}
\usepackage{caption}
\usepackage{subcaption}
\usepackage{tikz}
\usepackage{pgfplots}
\usepackage{pgfplotstable}
\usepackage{hyperref}

\pgfplotsset{compat=1.18}
\usepgfplotslibrary{groupplots}
\usetikzlibrary{arrows.meta,positioning,calc,fit,shapes.geometric}

\geometry{
  top=2.5cm,
  bottom=2.5cm,
  left=3cm,
  right=3cm
}

\onehalfspacing
\sloppy

\sisetup{
  detect-all,
  per-mode=symbol,
  group-separator={\,},
  output-decimal-marker={.}
}

\hypersetup{
  colorlinks=true,
  linkcolor=black,
  citecolor=black,
  urlcolor=blue
}

\definecolor{ctblue}{HTML}{1F77B4}
\definecolor{v2red}{HTML}{D62728}
\definecolor{fgmgreen}{HTML}{2CA02C}
\definecolor{gridgray}{HTML}{666666}
\definecolor{accentgold}{HTML}{B8860B}
\definecolor{softgray}{HTML}{F2F3F5}

\newcolumntype{L}[1]{>{\raggedright\arraybackslash}p{#1}}

\pgfplotsset{
  every axis/.append style={
    width=\textwidth,
    height=6.2cm,
    grid=both,
    grid style={line width=.1pt, draw=gray!20},
    major grid style={line width=.2pt, draw=gray!35},
    tick label style={font=\small},
    label style={font=\small},
    legend style={font=\small, draw=none, fill=white, fill opacity=0.85, text opacity=1},
    line width=1.1pt
  }
}

% ------------------------------------------------------------
% Datos usados por las figuras PGFPlots
% ------------------------------------------------------------
\pgfplotstableread{
idx SuCT SuV2 errpct tCT tV2 nCT nV2 speed
1 0.41898 0.42553 1.56 29.80 31.27 34 36 0.95
2 0.29214 0.29886 2.30 10.93 19.71 36 37 0.55
3 0.12836 0.11751 8.46 10.15 15.96 31 42 0.64
4 0.36139 0.38268 5.89 8.12 2.94 37 37 2.76
5 0.99143 0.98840 0.31 6.58 6.61 39 41 1.00
6 2.78845 2.52952 9.29 3.15 0.52 34 42 6.09
7 0.21001 0.20043 4.57 15.14 18.20 39 48 0.83
8 2.56077 2.31279 9.68 7.94 20.58 34 36 0.39
}\sweeptable

\pgfplotstableread{
phi SuCT SuV2 errpct
0.6 0.12836 0.11751 8.46
0.8 0.29214 0.29886 2.30
1.0 0.41898 0.42553 1.56
1.2 0.36139 0.38268 5.89
}\phitable

\pgfplotstableread{
phi Z Su npts time wnorm
1.0 0.05003 0.42553 36 28.29 0.225
1.1 0.05524 0.43306 36 2.78 0.018
1.2 0.06040 0.38559 36 1.22 0.458
1.3 0.06550 0.27246 38 2.58 0.493
1.4 0.07055 0.15553 41 2.91 0.811
1.5 0.07554 0.10706 43 3.51 0.125
}\fgmtable

\pgfplotstableread{
zmm T u
0.000 300.00 0.41898
39.000 300.00 0.41900
40.875 300.34 0.41967
41.625 337.00 0.47312
41.906 567.00 0.80427
42.141 1266.86 1.82660
42.234 1612.38 2.33806
42.375 1841.58 2.66721
42.469 1883.67 2.71815
43.125 2007.29 2.86856
44.250 2089.12 2.97023
48.000 2172.45 3.07080
60.000 2214.84 3.12078
120.000 2225.82 3.13359
}\profilect

\pgfplotstableread{
zmm T u
0.000 300.00 0.42553
38.686 300.01 0.42555
40.364 300.52 0.42654
41.036 344.47 0.49147
41.371 781.38 1.13270
41.475 1145.28 1.67338
41.604 1641.47 2.41829
41.682 1798.78 2.65050
41.889 1903.59 2.78413
42.200 1971.82 2.86793
43.857 2103.42 3.03402
48.000 2182.66 3.13079
60.000 2219.98 3.17542
120.000 2229.44 3.18661
}\profilev

\pgfplotstableread{
pass n Finf
0 26 3.415
1 30 2.410
2 34 0.000082
3 36 4.233
4 36 4.233
}\refinetable

\pgfplotstableread{
idx time
1 39.99
2 37.07
3 36.38
4 29.38
5 29.10
6 5.98
7 2.24
}\profiletime

% ------------------------------------------------------------
% Metadatos
% ------------------------------------------------------------
\title{\textbf{Resolución numérica de una llama libre premezclada unidimensional}\\
\large Desarrollo, validación frente a Cantera y generación de tablas FGM}
\author{Kevin Vidal}
\date{Julio de 2026}

\begin{document}
\pagenumbering{gobble}

\begin{titlepage}
  \centering
  \vspace*{1.5cm}
  {\Large Trabajo Fin de Máster\par}
  \vspace{1.4cm}
  {\huge\bfseries Resolución numérica de una llama libre premezclada unidimensional\par}
  \vspace{0.6cm}
  {\Large Desarrollo de solver propio, comparación con Cantera y generación de tablas FGM\par}
  \vfill
  {\Large Kevin Vidal\par}
  \vspace{0.4cm}
  {\large Julio de 2026\par}
  \vfill
  \resizebox{0.92\textwidth}{!}{%
  \begin{tikzpicture}[node distance=1.6cm, >=Latex]
    \tikzstyle{box}=[draw, rounded corners=2pt, minimum width=3.5cm, minimum height=1.05cm, align=center, fill=softgray]
    \node[box] (phys) {Modelo físico\\llama libre 1D};
    \node[box, right=of phys] (num) {Solver propio\\residual + Newton + malla};
    \node[box, right=of num] (val) {Validación\\Cantera};
    \node[box, below=of num] (fgm) {Aplicación\\tabla FGM $(Z,c)$};
    \draw[->, thick] (phys) -- (num);
    \draw[->, thick] (num) -- (val);
    \draw[->, thick] (num) -- (fgm);
  \end{tikzpicture}
  }
\end{titlepage}
\clearpage
\pagenumbering{roman}

\tableofcontents
\clearpage
\pagenumbering{arabic}

% ============================================================
\section{Introducción}
% ============================================================

La combustión premezclada aparece en turbinas de gas, motores, quemadores industriales y estrategias modernas de combustión pobre. En todos esos casos, la pregunta física central es cómo una mezcla fresca de combustible y oxidante se transforma en productos calientes mediante un frente de reacción muy delgado. La resolución directa de una llama tridimensional turbulenta con química detallada sigue siendo costosa para diseño e iteración de ingeniería; por ello, muchas metodologías de combustión computacional se apoyan en problemas canónicos unidimensionales que concentran la información termoquímica esencial.

El problema tratado en este Trabajo Fin de Máster es la llama libre premezclada, plana, laminar, estacionaria y unidimensional. Aunque parece un caso elemental, contiene las dificultades que hacen exigente la simulación de combustión: acoplamiento fuerte entre transporte y cinética, rigidez química, condiciones de contorno en dos extremos, adaptación de malla y una velocidad de propagación $S_u$ que no se impone como dato sino que emerge como valor propio del sistema.

El software de referencia para este problema es Cantera, que implementa llamas unidimensionales, modelos termoquímicos, cinética, transporte, Newton amortiguado, pseudo-transitorio y refinamiento de malla \cite{Cantera2023,CanteraDocs31}. En este trabajo Cantera se usa como competidor técnico y como patrón de validación, pero la contribución no consiste en llamar a Cantera: se construyó un solver propio, localizado en el directorio \texttt{V2/}, que reproduce la ruta de cálculo de una \texttt{FreeFlame} 1D y permite abrir, modificar y extender cada bloque numérico.

La originalidad del trabajo está en tres niveles. Primero, se formuló un residual propio para continuidad, energía y especies, con layout por punto y flujos difusivos corregidos. Segundo, se implementó una infraestructura numérica completa: Jacobiano local, alternativa block-tridiagonal, reutilización de factorizaciones, Newton amortiguado, avance pseudo-transitorio y refinamiento adaptativo. Tercero, se conectó el solver a una aplicación de mayor valor: generación de tablas \textit{Flamelet-Generated Manifold} (FGM) en el espacio $(Z,c)$, con cache de semillas, continuación entre equivalence ratios y comparación directa con barridos Cantera.

La memoria se organiza como sigue. La sección de objetivos fija el alcance. El marco teórico presenta el modelo de llama libre 1D, el transporte, la cinética y la idea FGM. La metodología documenta el código realmente implementado y contrasta cada decisión con Cantera. Los resultados muestran validación, tiempos, perfiles y generación de tablas. La discusión interpreta las diferencias, y las conclusiones recogen el aporte y las rutas de mejora.

% ============================================================
\section{Objetivos}
% ============================================================

\subsection{Objetivo general}

Desarrollar, documentar y validar un solver propio para la llama libre premezclada unidimensional, capaz de calcular perfiles de velocidad, temperatura, especies y velocidad laminar de llama, y demostrar su utilidad como base para generar tablas FGM comparables con las obtenidas mediante Cantera.

\subsection{Objetivos específicos}

\begin{enumerate}
  \item Formular el problema estacionario de una llama libre premezclada 1D como sistema acoplado de continuidad, energía y especies, identificando la velocidad de llama como valor propio.
  \item Implementar un residual discreto propio con transporte \textit{mixture-averaged}, corrección de flujo difusivo y tratamiento de condiciones de contorno de entrada, salida y anclaje térmico.
  \item Construir una ruta numérica robusta con Newton amortiguado, pseudo-transitorio, BDF1/BDF2, reutilización del Jacobiano y refinamiento adaptativo de malla.
  \item Implementar y comparar diferentes estrategias de Jacobiano y solución lineal, incluyendo una ruta local tipo Cantera y una ruta block-tridiagonal propia.
  \item Incorporar una ruta termoquímica nativa basada en polinomios NASA-7, cinética Arrhenius con reacciones de tres cuerpos/falloff y transporte \textit{mixture-averaged}.
  \item Validar el solver frente a Cantera para casos de CH$_4$/aire e H$_2$/aire, variando equivalence ratio, temperatura de entrada y presión.
  \item Generar tablas FGM en el espacio $(Z,c)$ mediante barridos de flamelets, variable de progreso, fracción de mezcla de Bilger, cache de semillas y continuación.
  \item Analizar precisión, coste computacional, sensibilidad de malla y limitaciones actuales del solver propio frente a Cantera.
  \item Presentar la metodología y los resultados con figuras generadas en LaTeX mediante \texttt{pgfplots} y TikZ, evitando depender de imágenes externas para la memoria.
\end{enumerate}

% ============================================================
\section{Marco teórico}
% ============================================================

\subsection{Llama libre premezclada 1D}

Una llama premezclada plana se define por una mezcla fresca de combustible y oxidante que entra al dominio con composición uniforme. En régimen laminar y estacionario, las variables dependen solamente de la coordenada axial $z\in[0,L]$. El frente de llama separa la zona de reactivos frescos de la zona de productos quemados, y su estructura resulta del equilibrio local entre difusión de calor, difusión de especies y reacción química.

La masa total cumple
\begin{equation}
  \frac{d}{dz}(\rho u)=0,
  \qquad
  \dot{m}=\rho u=\rho_u S_u,
  \label{eq:continuidad}
\end{equation}
donde $\rho$ es la densidad, $u$ la velocidad axial, $\dot{m}$ el flujo másico y $S_u$ la velocidad laminar de llama respecto a la mezcla fresca.

La energía, en forma estacionaria y a presión prácticamente constante, se escribe como
\begin{equation}
  \rho c_p u \frac{dT}{dz}
  =
  \frac{d}{dz}\left(\lambda \frac{dT}{dz}\right)
  - \sum_{k=1}^{K} h_k \dot{\omega}_k
  - \sum_{k=1}^{K} J_k \frac{dh_k}{dz},
  \label{eq:energia}
\end{equation}
donde $T$ es la temperatura, $c_p$ el calor específico de mezcla, $\lambda$ la conductividad, $h_k$ la entalpía específica, $J_k$ el flujo difusivo y $\dot{\omega}_k$ la producción másica de la especie $k$.

Para cada especie,
\begin{equation}
  \rho u \frac{dY_k}{dz}
  =
  -\frac{dJ_k}{dz}+\dot{\omega}_k,
  \qquad k=1,\ldots,K,
  \label{eq:especies}
\end{equation}
con $Y_k$ la fracción másica. Las ecuaciones \eqref{eq:continuidad}-\eqref{eq:especies} forman un sistema rígido de ecuaciones diferenciales y algebraicas acopladas. La rigidez se origina en que las escalas químicas pueden ser muchos órdenes de magnitud menores que las escalas convectivas y difusivas \cite{PoinsotVeynante2005,Kee2017}.

\subsection{Valor propio y anclaje térmico}

En una llama libre, $S_u$ no se prescribe. El sistema admite una degeneración por traslación: si una solución es válida, la misma solución desplazada dentro del dominio también lo es. Para cerrar el problema se fija una temperatura en un punto interior,
\begin{equation}
  T(z_\mathrm{fix})=T^\ast,
  \label{eq:anchor}
\end{equation}
lo que determina la posición del frente y permite que $\dot{m}$, y por tanto $S_u$, se calcule como incógnita del sistema. Cantera aplica esta idea en su configuración de flujo libre, usando un punto interno de temperatura fija para determinar el flujo másico de entrada \cite{CanteraDocs31}.

\subsection{Transporte \textit{mixture-averaged}}

El transporte difusivo se modela con coeficientes \textit{mixture-averaged}. En una cara de celda $j+1/2$, el flujo preliminar de especie se calcula, usando gradiente molar, como
\begin{equation}
  J_k^\ast
  =
  -\rho
  \frac{W_k}{W_\mathrm{mix}}
  D_k
  \frac{\partial X_k}{\partial z},
  \label{eq:flux_raw}
\end{equation}
donde $W_k$ es el peso molecular, $W_\mathrm{mix}$ el peso molecular de mezcla, $D_k$ el coeficiente de difusión de la especie en la mezcla y $X_k$ la fracción molar. Como la suma de flujos difusivos debe conservar masa, se aplica la corrección
\begin{equation}
  J_k = J_k^\ast - Y_k\sum_{\ell=1}^{K}J_\ell^\ast.
  \label{eq:flux_corregido}
\end{equation}
Esta corrección es una parte clave del solver propio porque evita que la discretización de especies introduzca una fuente artificial de masa.

\subsection{Cinética y propiedades termodinámicas}

El mecanismo de validación principal es GRI-Mech 3.0, que contiene 53 especies y 325 reacciones elementales para combustión de gas natural \cite{Smith1999}. Las constantes de velocidad directas se expresan con la forma de Arrhenius modificada,
\begin{equation}
  k_f = A T^b \exp\left(-\frac{E_a}{RT}\right),
  \label{eq:arrhenius}
\end{equation}
y las tasas inversas se obtienen por equilibrio químico cuando la reacción es reversible.

Las propiedades termodinámicas se calculan con polinomios NASA-7 \cite{McBride2002}. Para una especie, en cada intervalo de temperatura,
\begin{equation}
  \frac{c_{p,k}^0}{R}
  =
  a_1+a_2T+a_3T^2+a_4T^3+a_5T^4,
  \label{eq:nasa_cp}
\end{equation}
mientras que la entalpía y entropía se obtienen por integración analítica. Esta estructura es conveniente para un solver propio porque transforma la evaluación termoquímica en operaciones vectorizables y diferenciables por perturbación local.

\subsection{Cantera como referencia}

Cantera resuelve flujos reactivos 1D con dominios compuestos, esquemas de diferencias finitas, Newton amortiguado, integración pseudo-transitoria y refinamiento de malla \cite{Cantera2023,CanteraDocs31}. En particular, \texttt{FreeFlame} representa una llama plana libre y permite resolver llamas premezcladas con refinamiento automático. En esta tesis Cantera cumple tres papeles:

\begin{itemize}
  \item referencia física y numérica para validar $S_u$, perfiles, número de puntos y comportamiento de malla;
  \item competidor de rendimiento, al medir tiempos de solución para los mismos casos;
  \item fuente de diseño algorítmico, porque varias ideas robustas de Cantera se reimplementan de forma explícita y editable en el solver propio.
\end{itemize}

La comparación no pretende afirmar que el código propio sustituye por completo a Cantera. La contribución es disponer de un núcleo abierto, controlable y extensible, diseñado para incorporar términos adicionales como pérdidas de energía, modelos de hollín o estrategias FGM específicas.

\subsection{Tablas FGM}

El método \textit{Flamelet-Generated Manifold} reduce el coste de simulaciones multidimensionales almacenando la termoquímica detallada en una tabla de baja dimensión construida con flamelets unidimensionales \cite{vanOijen2000}. En este trabajo se usa una tabla en el espacio $(Z,c)$, donde $Z$ es la fracción de mezcla y $c$ una variable de progreso.

La fracción de mezcla de Bilger combina los elementos C, H y O en una cantidad conservada \cite{Bilger1989}. En Cantera también se documenta la opción de calcular $Z$ mediante la mezcla de Bilger \cite{CanteraDocs31}. Para el caso CH$_4$/aire, $Z$ se relaciona de forma monótona con $\phi$, por lo que un barrido en equivalence ratio puede almacenarse como un eje de mezcla.

La variable de progreso usada es
\begin{equation}
  c=\frac{\beta-\beta_u}{\beta_b-\beta_u},
  \qquad
  \beta=Y_{\mathrm{CO_2}}+Y_{\mathrm{H_2O}}+Y_{\mathrm{CO}}+0.5Y_{\mathrm{H_2}},
  \label{eq:progress}
\end{equation}
donde $u$ y $b$ indican mezcla fresca y productos quemados. Esta definición concentra la evolución química principal y permite interpolar temperatura, especies, calor liberado y otras magnitudes sobre una malla adaptativa en $c$.

% ============================================================
\section{Metodología}
% ============================================================

\subsection{Arquitectura implementada}

El trabajo se implementó en una arquitectura modular. El flujo principal reside en \texttt{V2/}; la generación FGM reside en \texttt{FGM/scripts/}. Los módulos principales son:

\begin{itemize}
  \item \texttt{problem.py}: define el caso físico, composición de entrada, equilibrio adiabático, malla inicial, anclaje térmico, límites de variables y tolerancias.
  \item \texttt{state.py}: empaqueta y desempaqueta el estado con layout por punto, $[u,T,Y_1,\ldots,Y_K]$ en cada nodo.
  \item \texttt{equations.py}: ensambla residual, flujos, condiciones de contorno, Jacobianos y solucionadores lineales.
  \item \texttt{solver.py}: contiene Newton amortiguado, pseudo-transitorio, expansión de dominio, refinamiento adaptativo y estrategia global tipo \texttt{auto}.
  \item \texttt{materiales/}: implementa termodinámica NASA-7, cinética, transporte y backend nativo sin depender de Cantera para cada evaluación local.
  \item \texttt{generate\_fgm\_tables\_jax.py}: genera tablas FGM usando el solver propio V2, cache de semillas, continuación y tabla $(Z,c)$.
  \item \texttt{generate\_fgm\_tables\_cantera.py}: genera la contraparte de referencia con Cantera.
\end{itemize}

\begin{figure}[H]
  \centering
  \resizebox{\textwidth}{!}{%
  \begin{tikzpicture}[
    node distance=1.05cm and 1.15cm,
    >=Latex,
    block/.style={draw, rounded corners=2pt, align=center, fill=softgray, minimum height=0.85cm, text width=3.0cm, font=\small},
    key/.style={draw, rounded corners=2pt, align=center, fill=white, minimum height=0.85cm, text width=3.2cm, font=\small}
  ]
    \node[block] (case) {Caso físico\\$T_u,p,\phi$, mecanismo};
    \node[block, right=of case] (prob) {\texttt{problem.py}\\malla, ancla, límites};
    \node[block, right=of prob] (res) {\texttt{equations.py}\\residual y Jacobiano};
    \node[block, below=of res] (solv) {\texttt{solver.py}\\Newton + transitorio};
    \node[block, left=of solv] (back) {Backend nativo\\NASA, cinética, transporte};
    \node[block, right=of solv] (out) {Perfiles\\$u,T,Y_k,S_u$};
    \node[key, below=of solv] (fgm) {FGM\\tabla $(Z,c)$};
    \node[key, above=of res] (ct) {Cantera\\referencia};
    \draw[->, thick] (case) -- (prob);
    \draw[->, thick] (prob) -- (res);
    \draw[->, thick] (back) -- (res);
    \draw[->, thick] (res) -- (solv);
    \draw[->, thick] (solv) -- (out);
    \draw[->, thick] (out) -- (fgm);
    \draw[->, thick, dashed] (ct) -- (res);
    \draw[->, thick, dashed] (ct) -- (out);
  \end{tikzpicture}
  }
  \caption{Arquitectura del solver propio y puntos de comparación frente a Cantera.}
  \label{fig:pipeline}
\end{figure}

\subsection{Discretización y layout del estado}

El dominio se discretiza en $N$ nodos no uniformes. El estado se almacena con layout por punto:
\begin{equation}
  \mathbf{x}
  =
  [u_0,T_0,Y_{1,0},\ldots,Y_{K,0},
   u_1,T_1,Y_{1,1},\ldots,Y_{K,1},\ldots]^\top.
  \label{eq:layout}
\end{equation}
Este diseño es deliberado: coincide con la forma conceptual de los dominios 1D de Cantera y hace que cada ecuación dependa solamente del nodo actual y sus vecinos inmediatos. Por tanto, el Jacobiano tiene estructura de banda con bloques vecinos, una propiedad explotada posteriormente.

La convección se aproxima con upwind de primer orden, la conducción y difusión con diferencias centradas en malla no uniforme y los flujos se evalúan en caras. En el extremo izquierdo se fija la temperatura de entrada y se impone el balance de especies con la composición fresca. En el extremo derecho se usa gradiente nulo o condición de flujo de especies tipo Cantera según la opción del caso. En el punto de anclaje se reemplaza la ecuación de continuidad por \eqref{eq:anchor}.

\subsection{Termoquímica nativa}

Una parte original del trabajo fue separar dos rutas de propiedades:

\begin{itemize}
  \item una ruta de referencia basada en Cantera, útil para comparación y depuración;
  \item una ruta nativa en \texttt{materiales/}, que evalúa termodinámica, transporte y cinética directamente desde los datos del mecanismo.
\end{itemize}

La ruta nativa implementa polinomios NASA-7, reacciones elementales, reversibles, de tres cuerpos y falloff tipo Lindemann/Troe. El transporte usa propiedades \textit{mixture-averaged}, reglas de mezcla y coeficientes polinomiales exportados, con kernels Numba en los bucles más costosos. Esta parte es importante porque convierte el TFM en una plataforma extensible: si se añaden pérdidas térmicas, especies de hollín o términos de radiación, no es necesario esperar a que Cantera exponga exactamente el modelo deseado.

\subsection{Jacobiano local y estructura block-tridiagonal}

El Jacobiano se construye por perturbación local. Para una columna asociada al nodo $j$, solamente se recalculan las filas de los nodos $j-1$, $j$ y $j+1$. Así se conserva la idea robusta de Cantera, pero con control explícito de la memoria y de los bloques.

Si $n_v=K+2$ es el número de variables por punto, el Jacobiano tiene la forma
\begin{equation}
  \mathbf{J}
  =
  \begin{bmatrix}
  D_0 & U_0 & 0 & \cdots & 0\\
  L_1 & D_1 & U_1 & \ddots & \vdots\\
  0 & L_2 & D_2 & \ddots & 0\\
  \vdots & \ddots & \ddots & \ddots & U_{N-2}\\
  0 & \cdots & 0 & L_{N-1} & D_{N-1}
  \end{bmatrix},
  \label{eq:blocktridiag}
\end{equation}
donde cada bloque es de tamaño $n_v\times n_v$. El modo \texttt{block\_tridiag} almacena directamente los bloques $L_j$, $D_j$ y $U_j$, evitando tratar la matriz como dispersa genérica cuando no es necesario.

\begin{figure}[H]
  \centering
  \begin{tikzpicture}[scale=0.55]
    \foreach \i in {0,...,6} {
      \draw[gray!35] (0,-\i) -- (7,-\i);
      \draw[gray!35] (\i,0) -- (\i,-7);
    }
    \foreach \i in {0,...,5} {
      \fill[ctblue!70] (\i,-\i) rectangle ++(0.82,-0.82);
    }
    \foreach \i in {0,...,4} {
      \fill[v2red!70] (\i+1,-\i) rectangle ++(0.82,-0.82);
      \fill[fgmgreen!70] (\i,-\i-1) rectangle ++(0.82,-0.82);
    }
    \node[font=\small] at (2.9,0.55) {bloques por nodo};
    \node[font=\small, ctblue] at (7.6,-2.0) {$D_j$};
    \node[font=\small, v2red] at (7.6,-2.7) {$U_j$};
    \node[font=\small, fgmgreen] at (7.6,-3.4) {$L_j$};
  \end{tikzpicture}
  \caption{Estructura block-tridiagonal inducida por la discretización 1D con acoplamiento local.}
  \label{fig:blocktridiag}
\end{figure}

\subsection{Newton amortiguado y pseudo-transitorio}

El sistema no lineal se resuelve con Newton amortiguado. En cada iteración se resuelve
\begin{equation}
  \mathbf{J}\Delta\mathbf{x}=-\mathbf{F}(\mathbf{x}),
  \label{eq:newton}
\end{equation}
se limita el paso para respetar cotas físicas de temperatura y fracciones másicas, y se aplica una búsqueda de amortiguación hasta que la norma ponderada del paso mejora. La norma usada es
\begin{equation}
  \|\Delta\mathbf{x}\|_w
  =
  \left[
  \frac{1}{Nn_v}
  \sum_i
  \left(
  \frac{\Delta x_i}{r_\mathrm{tol}\,\overline{|x_i|}+a_\mathrm{tol}}
  \right)^2
  \right]^{1/2}.
  \label{eq:weightednorm}
\end{equation}

Cuando Newton estacionario falla, el solver alterna pasos pseudo-transitorios. Con BDF1,
\begin{equation}
  \mathbf{F}_t(\mathbf{x})
  =
  \mathbf{F}_{ss}(\mathbf{x})
  -
  \mathbf{m}\odot
  \frac{\mathbf{x}-\mathbf{x}^n}{\Delta t},
  \label{eq:bdf1}
\end{equation}
y con BDF2,
\begin{equation}
  \mathbf{F}_t(\mathbf{x})
  =
  \mathbf{F}_{ss}(\mathbf{x})
  -
  \mathbf{m}\odot
  \frac{3\mathbf{x}-4\mathbf{x}^n+\mathbf{x}^{n-1}}{2\Delta t}.
  \label{eq:bdf2}
\end{equation}
La máscara $\mathbf{m}$ activa solamente variables diferenciales, temperatura y especies, y deja la velocidad como variable algebraica. El Jacobiano estacionario se reutiliza y solo se actualiza su diagonal transitoria cuando cambia $\Delta t$.

\subsection{Expansión de dominio y refinamiento}

La malla inicial puede ser la canónica de Cantera o una malla propia con concentración gaussiana en la región esperada de llama. Luego se aplica refinamiento adaptativo con criterios de razón, pendiente y curvatura:
\begin{align}
  \frac{h_{j+1}}{h_j} &> r,
  \label{eq:ratio}\\
  \frac{|\phi_{j+1}-\phi_j|}{\max |\Delta\phi|} &> \sigma_s,
  \label{eq:slope}\\
  \frac{|s_{j+1}-s_j|}{\max |\Delta s|} &> \sigma_c.
  \label{eq:curve}
\end{align}
Los valores base usados en validación son $r=10$, $\sigma_s=0.8$, $\sigma_c=0.8$ y poda desactivada, alineados con el modo de referencia. El solver también expande el dominio si detecta que el frente no queda suficientemente contenido.

\subsection{Generación FGM}

La generación FGM usa el solver propio de forma repetida. Para cada $\phi$ o $Z$ objetivo:

\begin{enumerate}
  \item se construye el caso de llama libre;
  \item se carga una semilla de cache si existe;
  \item se aplica continuación desde el flamelet anterior, con predictor secante amortiguado;
  \item se resuelve la llama con el solver V2;
  \item se reparametriza el perfil por $c$;
  \item se interpola en la malla global $(Z,c)$;
  \item se guardan \texttt{fgm\_table.npz}, \texttt{summary\_phi.csv}, perfiles crudos y metadatos.
\end{enumerate}

La originalidad de esta parte no es la idea FGM en sí, que está establecida en la literatura \cite{vanOijen2000}, sino su integración con un solver propio y con mecanismos prácticos de producción: cache de semillas, continuación, inversión $Z\leftrightarrow\phi$, malla adaptativa en $c$, y exportación posterior a formatos de flamelets.

\subsection{Diseño de comparación frente a Cantera}

La comparación se hizo en tres niveles:

\begin{itemize}
  \item \textbf{Caso base}: CH$_4$/aire, $\phi=1.0$, $T_u=\SI{300}{K}$, $p=\SI{1}{atm}$, GRI-Mech 3.0, transporte \textit{mixture-averaged}.
  \item \textbf{Barrido de validación}: variación de $\phi$, temperatura de entrada, presión y combustible, comparando $S_u$, tiempo y puntos de malla.
  \item \textbf{Barrido FGM}: generación de flamelets ricos CH$_4$/aire para $\phi=1.0$ a $1.5$, evaluando convergencia, tiempo, puntos y continuidad de $S_u$.
\end{itemize}

Los resultados numéricos proceden de las corridas de comparación de V2 y de los barridos FGM guardados en el proyecto. Las figuras de esta memoria se generan directamente con \texttt{pgfplots} a partir de datos tabulados en el propio archivo \LaTeX.

% ============================================================
\section{Resultados}
% ============================================================

\subsection{\texorpdfstring{Caso base CH$_4$/aire}{Caso base CH4/aire}}

Para la llama estequiométrica CH$_4$/aire a \SI{300}{K} y \SI{1}{atm}, Cantera predice $S_u=\SI{0.41898}{m/s}$ con 34 puntos, mientras que el solver propio obtiene $S_u=\SI{0.42553}{m/s}$ con 36 puntos. El error relativo en velocidad de llama es \SI{1.56}{\percent}. Esta coincidencia es relevante porque $S_u$ depende de la integración completa de transporte, cinética, energía y condiciones de contorno.

\begin{figure}[H]
  \centering
  \begin{tikzpicture}
    \begin{groupplot}[
      group style={group size=1 by 2, vertical sep=1.25cm},
      width=0.93\textwidth,
      height=5.3cm,
      xmin=38, xmax=46,
      legend columns=2,
      legend pos=north west
    ]
      \nextgroupplot[
        ylabel={$T$ [K]},
        xlabel={$z$ [mm]},
        ymin=250, ymax=2350
      ]
      \addplot[color=ctblue, mark=*] table[x=zmm,y=T] {\profilect};
      \addlegendentry{Cantera}
      \addplot[color=v2red, mark=square*] table[x=zmm,y=T] {\profilev};
      \addlegendentry{V2 propio}
      \nextgroupplot[
        ylabel={$u$ [m/s]},
        xlabel={$z$ [mm]},
        ymin=0, ymax=3.4
      ]
      \addplot[color=ctblue, mark=*] table[x=zmm,y=u] {\profilect};
      \addlegendentry{Cantera}
      \addplot[color=v2red, mark=square*] table[x=zmm,y=u] {\profilev};
      \addlegendentry{V2 propio}
    \end{groupplot}
  \end{tikzpicture}
  \caption{Perfiles de temperatura y velocidad del caso base CH$_4$/aire, $\phi=1.0$. La posición del frente presenta un pequeño desplazamiento de anclaje, pero la estructura y los niveles quemados son coherentes.}
  \label{fig:profiles_base}
\end{figure}

La comparación directa de perfiles en coordenada absoluta puede sobrestimar diferencias de temperatura cuando el frente está desplazado unas décimas de milímetro. En perfiles alineados por el máximo de liberación de calor, las corridas guardadas muestran desplazamientos de \SI{0.651}{mm} para $\phi=1.0$ y \SI{0.677}{mm} para $\phi=1.05$, con errores relativos de $S_u$ de \SI{1.56}{\percent} y \SI{1.49}{\percent}, respectivamente.

\subsection{Barrido de validación}

La Tabla~\ref{tab:sweep} resume ocho casos de validación. Se incluyen variaciones de equivalence ratio, precalentamiento, presión y combustible. El solver propio converge en todos los casos guardados, con errores de $S_u$ entre \SI{0.31}{\percent} y \SI{9.68}{\percent}.

\begin{table}[H]
  \centering
  \caption{Comparación de velocidad de llama y coste frente a Cantera.}
  \label{tab:sweep}
  \small
  \begin{tabular}{lrrrrr}
    \toprule
    Caso & $S_{u,\mathrm{CT}}$ & $S_{u,\mathrm{V2}}$ & Error [\%] & $t_\mathrm{CT}$ [s] & $t_\mathrm{V2}$ [s]\\
    \midrule
    CH$_4$ $\phi=1.0$ & 0.41898 & 0.42553 & 1.56 & 29.80 & 31.27\\
    CH$_4$ $\phi=0.8$ & 0.29214 & 0.29886 & 2.30 & 10.93 & 19.71\\
    CH$_4$ $\phi=0.6$ & 0.12836 & 0.11751 & 8.46 & 10.15 & 15.96\\
    CH$_4$ $\phi=1.2$ & 0.36139 & 0.38268 & 5.89 & 8.12 & 2.94\\
    CH$_4$ $T_u=\SI{500}{K}$ & 0.99143 & 0.98840 & 0.31 & 6.58 & 6.61\\
    CH$_4$ $T_u=\SI{800}{K}$ & 2.78845 & 2.52952 & 9.29 & 3.15 & 0.52\\
    CH$_4$ $p=\SI{5}{atm}$ & 0.21001 & 0.20043 & 4.57 & 15.14 & 18.20\\
    H$_2$ $\phi=1.0$ & 2.56077 & 2.31279 & 9.68 & 7.94 & 20.58\\
    \bottomrule
  \end{tabular}
\end{table}

\begin{figure}[H]
  \centering
  \begin{tikzpicture}
    \begin{axis}[
      ybar,
      bar width=5pt,
      ylabel={$S_u$ [m/s]},
      xlabel={Caso},
      symbolic x coords={c1,c2,c3,c4,c5,c6,c7,c8},
      xtick=data,
      xticklabels={CH$_4$ $\phi$1.0,CH$_4$ $\phi$0.8,CH$_4$ $\phi$0.6,CH$_4$ $\phi$1.2,CH$_4$ 500K,CH$_4$ 800K,CH$_4$ 5atm,H$_2$ $\phi$1.0},
      x tick label style={rotate=40, anchor=east, font=\scriptsize},
      ymin=0,
      legend pos=north west
    ]
      \addplot[fill=ctblue!70, draw=ctblue] coordinates {
        (c1,0.41898) (c2,0.29214) (c3,0.12836) (c4,0.36139)
        (c5,0.99143) (c6,2.78845) (c7,0.21001) (c8,2.56077)
      };
      \addlegendentry{Cantera}
      \addplot[fill=v2red!70, draw=v2red] coordinates {
        (c1,0.42553) (c2,0.29886) (c3,0.11751) (c4,0.38268)
        (c5,0.98840) (c6,2.52952) (c7,0.20043) (c8,2.31279)
      };
      \addlegendentry{V2 propio}
    \end{axis}
  \end{tikzpicture}
  \caption{Velocidad laminar de llama calculada por Cantera y por el solver propio en el barrido de validación.}
  \label{fig:su_cases}
\end{figure}

El barrido de CH$_4$/aire con $\phi=0.6$ a $1.2$ reproduce la tendencia física: $S_u$ aumenta desde mezcla pobre hasta la zona cercana a estequiometría y luego disminuye para mezcla rica. El error más bajo se obtiene en $\phi=1.0$ y $\phi=0.8$; el error aumenta en $\phi=0.6$ y $\phi=1.2$, donde la sensibilidad a malla, química y frontera de salida es mayor.

\begin{figure}[H]
  \centering
  \begin{tikzpicture}
    \begin{groupplot}[
      group style={group size=1 by 2, vertical sep=1.2cm},
      width=\textwidth,
      height=5.1cm,
      xmin=0.55, xmax=1.25,
      xtick={0.6,0.8,1.0,1.2}
    ]
      \nextgroupplot[
        ylabel={$S_u$ [m/s]},
        xlabel={$\phi$},
        legend pos=north west
      ]
      \addplot[color=ctblue, mark=*] table[x=phi,y=SuCT] {\phitable};
      \addlegendentry{Cantera}
      \addplot[color=v2red, mark=square*] table[x=phi,y=SuV2] {\phitable};
      \addlegendentry{V2 propio}
      \nextgroupplot[
        ylabel={Error relativo [\%]},
        xlabel={$\phi$},
        ymin=0
      ]
      \addplot[color=accentgold, mark=triangle*] table[x=phi,y=errpct] {\phitable};
    \end{groupplot}
  \end{tikzpicture}
  \caption{Validación del barrido CH$_4$/aire frente a Cantera.}
  \label{fig:phi_sweep}
\end{figure}

\subsection{Refinamiento de malla}

En el caso base, el solver expandió el dominio y refinó la malla hasta 36 puntos. La secuencia de refinamiento registrada fue $22\rightarrow 26\rightarrow 30\rightarrow 34\rightarrow 36$ puntos, con aceptación final al no requerir nuevos puntos. Esta traza confirma que el refinador no es decorativo: participa en la convergencia del perfil y en la localización del frente.

\begin{figure}[H]
  \centering
  \begin{tikzpicture}
    \begin{groupplot}[
      group style={group size=1 by 2, vertical sep=1.15cm},
      width=0.93\textwidth,
      height=4.8cm,
      xmin=-0.2, xmax=4.2,
      xtick={0,1,2,3,4}
    ]
      \nextgroupplot[
        ylabel={Puntos},
        xlabel={Pase de refinamiento},
        ymin=20, ymax=40
      ]
      \addplot[color=ctblue, mark=*] table[x=pass,y=n] {\refinetable};
      \nextgroupplot[
        ylabel={$||F||_\infty$},
        xlabel={Pase de refinamiento},
        ymode=log,
        ymin=1e-5, ymax=10
      ]
      \addplot[color=v2red, mark=square*] table[x=pass,y=Finf] {\refinetable};
    \end{groupplot}
  \end{tikzpicture}
  \caption{Evolución de malla y residual infinito durante el refinamiento del caso base.}
  \label{fig:refinement}
\end{figure}

\subsection{Coste computacional}

El perfil de tiempo instrumentado del caso base muestra que el coste dominante se concentra en construir y usar el modelo lineal: Jacobiano, residuales locales, residuales completos y amortiguación de Newton. Las categorías no deben sumarse de forma estricta porque algunas mediciones están anidadas, pero sí identifican el cuello de botella.

\begin{figure}[H]
  \centering
  \begin{tikzpicture}
    \begin{axis}[
      ybar,
      bar width=10pt,
      ylabel={Tiempo instrumentado [s]},
      symbolic x coords={linear,jac,full,damp,local,fact,solve},
      xtick=data,
      xticklabels={modelo lineal,Jacobiano,residual completo,damping,residual local,LU,solve},
      x tick label style={rotate=35, anchor=east, font=\scriptsize},
      ymin=0
    ]
      \addplot[fill=gridgray!70, draw=gridgray] coordinates {
        (linear,39.99)
        (jac,37.07)
        (full,36.38)
        (damp,29.38)
        (local,29.10)
        (fact,5.98)
        (solve,2.24)
      };
    \end{axis}
  \end{tikzpicture}
  \caption{Perfil de coste instrumentado para el caso base. El foco de optimización queda en Jacobiano, residual y amortiguación.}
  \label{fig:profile_time}
\end{figure}

\subsection{Generación de tabla FGM}

El barrido FGM propio para CH$_4$/aire rico resolvió seis flamelets entre $\phi=1.0$ y $\phi=1.5$, todos aceptados con criterio tipo Cantera basado en norma ponderada del paso y guardia de residual. El primer flamelet costó \SI{28.29}{s}; los siguientes bajaron a entre \SI{1.22}{s} y \SI{3.51}{s}, mostrando el efecto de cache y continuación.

\begin{table}[H]
  \centering
  \caption{Barrido FGM propio en espacio $(Z,c)$ para CH$_4$/aire.}
  \label{tab:fgm}
  \small
  \begin{tabular}{rrrrrr}
    \toprule
    $\phi$ & $Z$ & $S_u$ [m/s] & Puntos & Tiempo [s] & Norma ponderada\\
    \midrule
    1.0 & 0.05003 & 0.42553 & 36 & 28.29 & 0.225\\
    1.1 & 0.05524 & 0.43306 & 36 & 2.78 & 0.018\\
    1.2 & 0.06040 & 0.38559 & 36 & 1.22 & 0.458\\
    1.3 & 0.06550 & 0.27246 & 38 & 2.58 & 0.493\\
    1.4 & 0.07055 & 0.15553 & 41 & 2.91 & 0.811\\
    1.5 & 0.07554 & 0.10706 & 43 & 3.51 & 0.125\\
    \bottomrule
  \end{tabular}
\end{table}

\begin{figure}[H]
  \centering
  \begin{tikzpicture}
    \begin{groupplot}[
      group style={group size=1 by 2, vertical sep=1.2cm},
      width=\textwidth,
      height=5.0cm,
      xmin=0.95, xmax=1.55,
      xtick={1.0,1.1,1.2,1.3,1.4,1.5}
    ]
      \nextgroupplot[
        ylabel={$S_u$ [m/s]},
        xlabel={$\phi$}
      ]
      \addplot[color=fgmgreen, mark=*] table[x=phi,y=Su] {\fgmtable};
      \nextgroupplot[
        ylabel={Tiempo [s]},
        xlabel={$\phi$},
        ymin=0
      ]
      \addplot[color=v2red, mark=square*] table[x=phi,y=time] {\fgmtable};
    \end{groupplot}
  \end{tikzpicture}
  \caption{Barrido FGM propio: velocidad de llama y tiempo por flamelet. El primer punto absorbe el coste de arranque; la continuación acelera los siguientes.}
  \label{fig:fgm_sweep}
\end{figure}

La tabla FGM final almacena ejes $Z$ y $c$, temperatura, velocidad, densidad, calor liberado y fracciones másicas. Esto convierte el solver 1D en un generador de datos para simulaciones multidimensionales, que era uno de los objetivos centrales del proyecto.

% ============================================================
\section{Discusión}
% ============================================================

\subsection{Qué se consiguió frente a Cantera}

Cantera sigue siendo más general y más maduro. Sin embargo, los resultados muestran que el solver propio ya reproduce la física principal de una llama libre premezclada y alcanza errores bajos en casos clave. El caso CH$_4$/aire estequiométrico queda a \SI{1.56}{\percent} de Cantera en $S_u$, y el caso precalentado a \SI{500}{K} queda a \SI{0.31}{\percent}. Estos dos puntos son importantes porque combinan química detallada, transporte, energía, adaptación de malla y cálculo de valor propio.

La aportación original aparece donde Cantera actúa como caja robusta y el solver propio como plataforma editable. En el código desarrollado se puede modificar directamente el residual, congelar o descongelar transporte en el Jacobiano, cambiar la condición de salida, añadir términos fuente, instrumentar tiempos y exportar datos FGM. Esta apertura es la razón técnica de construir una herramienta propia incluso si Cantera existe.

\subsection{Diferencias y limitaciones}

Los errores mayores aparecen en CH$_4$ pobre, CH$_4$ muy precalentado e H$_2$/aire. Estos casos son numéricamente más sensibles: las capas son más delgadas o las difusividades alteran de forma fuerte la estructura del frente. En H$_2$, además, la ausencia de modelos multicomponente/Soret en la ruta base puede afectar más que en CH$_4$.

La comparación de perfiles también muestra que una diferencia pequeña de posición del frente genera errores puntuales grandes en temperatura si se comparan coordenadas absolutas. Esto no invalida el $S_u$, pero exige comparar perfiles alineados por un marcador físico, como el máximo de liberación de calor, para separar error de estructura y error de fase.

En rendimiento, el solver propio todavía no es uniformemente más rápido. En algunos casos iguala o supera a Cantera, por ejemplo CH$_4$ $\phi=1.2$ y CH$_4$ a \SI{800}{K}; en otros es más lento. El perfil indica que la ruta de mejora no está en cambios cosméticos, sino en reducir coste de Jacobiano, residual local y amortiguación. La implementación block-tridiagonal, el cache de semillas y la continuación ya apuntan en esa dirección.

\subsection{Valor de la ruta FGM}

La generación FGM es una demostración de utilidad más allá del caso 1D aislado. Cantera puede generar flamelets, pero el flujo propio agrega control sobre:

\begin{itemize}
  \item cómo se escoge y amortigua la semilla entre flamelets;
  \item cómo se acepta una solución cuando el criterio ponderado tipo Cantera es razonable aunque el residual bruto sea grande;
  \item qué variables se guardan en la tabla y cómo se refina el eje $c$;
  \item cómo se exporta la información para otros solvers.
\end{itemize}

Por tanto, el resultado innovador no es solamente resolver una llama: es haber construido una cadena completa que va desde ecuaciones diferenciales hasta una tabla termoquímica reutilizable.

\subsection{Originalidad por etapas}

\begin{table}[H]
  \centering
  \caption{Comparación conceptual entre Cantera y la implementación propia.}
  \label{tab:originalidad}
  \small
  \begin{tabular}{L{0.25\textwidth}L{0.31\textwidth}L{0.35\textwidth}}
    \toprule
    Etapa & Cantera & Implementación propia\\
    \midrule
    Formulación & Modelo 1D robusto y general & Residual explícito en \texttt{equations.py}, ajustable para nuevos términos\\
    Termoquímica & Librería integrada & Ruta nativa NASA-7, cinética y transporte en \texttt{materiales/}\\
    Jacobiano & Local y disperso & Local, batched, banded y block-tridiagonal configurable\\
    Newton & Híbrido con pseudo-transitorio & Reimplementación inspectable con edad de Jacobiano, BDF1/BDF2 y métricas\\
    Malla & Refinador automático & Refinador propio con perfiles seleccionables y trazas de refinamiento\\
    FGM & Puede resolver flamelets & Pipeline propio con $Z$, $c$, cache, continuación y exportación\\
    \bottomrule
  \end{tabular}
\end{table}

Esta tabla resume la tesis técnica del trabajo: Cantera es el estándar de comparación, pero el valor del TFM está en haber hecho explícito, modificable y medible el camino numérico completo.

% ============================================================
\section{Conclusiones}
% ============================================================

Se desarrolló un solver propio para la llama libre premezclada unidimensional, incluyendo formulación de residual, transporte corregido, química detallada, Newton amortiguado, pseudo-transitorio, refinamiento adaptativo y cálculo de $S_u$ como valor propio. El solver reproduce el comportamiento de Cantera en el caso base CH$_4$/aire con un error de \SI{1.56}{\percent} en velocidad de llama y alcanza errores menores al \SI{1}{\percent} en el caso precalentado a \SI{500}{K}.

La metodología implementada no es un envoltorio de Cantera. Se construyeron módulos propios para estado, residual, Jacobiano, solución no lineal, refinamiento y termoquímica nativa. Cantera se usó como referencia de validación y como competidor técnico, lo que permitió cuantificar precisión, coste y limitaciones.

La aplicación FGM confirma que el solver puede integrarse en un flujo de trabajo de combustión computacional. Se generó una tabla $(Z,c)$ con seis flamelets ricos, todos aceptados, y se observó que la continuación y el cache reducen fuertemente el coste después del primer punto.

Las limitaciones actuales se concentran en casos más sensibles, especialmente H$_2$/aire, mezclas muy pobres o temperaturas altas, y en el coste de Jacobiano/residual. Las extensiones prioritarias son: transporte multicomponente y Soret, términos de pérdida de energía, comparación de perfiles alineados por liberación de calor, Jacobianos semi-analíticos, paralelización de barridos FGM desde cache y validación con mecanismos reducidos.

En conjunto, el trabajo deja una base innovadora y extensible: un solver 1D entendible desde las ecuaciones hasta la tabla FGM, validado frente a Cantera y preparado para incorporar modelos que Cantera no expone de forma directa en el flujo deseado.

\clearpage
\phantomsection
\addcontentsline{toc}{section}{Referencias}
\begin{thebibliography}{99}

\bibitem{Cantera2023}
D. G. Goodwin, H. K. Moffat, I. Schoegl, R. L. Speth y B. W. Weber.
\newblock \textit{Cantera: An Object-oriented Software Toolkit for Chemical Kinetics, Thermodynamics, and Transport Processes}.
\newblock Zenodo, 2023. DOI: \href{https://doi.org/10.5281/zenodo.8137090}{10.5281/zenodo.8137090}.

\bibitem{CanteraDocs31}
Cantera Developers.
\newblock \textit{One-dimensional Reacting Flows}. Cantera 3.1.0 documentation.
\newblock \url{https://cantera.org/3.1/python/onedim.html}.

\bibitem{Kee2017}
R. J. Kee, M. E. Coltrin, P. Glarborg y H. Zhu.
\newblock \textit{Chemically Reacting Flow: Theory, Modeling, and Simulation}.
\newblock 2nd ed., Wiley, 2017.

\bibitem{PoinsotVeynante2005}
T. Poinsot y D. Veynante.
\newblock \textit{Theoretical and Numerical Combustion}.
\newblock 2nd ed., R.T. Edwards, 2005.

\bibitem{Smooke1991}
M. D. Smooke, editor.
\newblock \textit{Reduced Kinetic Mechanisms and Asymptotic Approximations for Methane-Air Flames}.
\newblock Lecture Notes in Physics, vol. 384, Springer, 1991.

\bibitem{Smith1999}
G. P. Smith, D. M. Golden, M. Frenklach y colaboradores.
\newblock \textit{GRI-Mech 3.0}.
\newblock Gas Research Institute, 1999. \url{https://combustion.berkeley.edu/gri-mech/}.

\bibitem{McBride2002}
B. J. McBride, M. J. Zehe y S. Gordon.
\newblock \textit{NASA Glenn Coefficients for Calculating Thermodynamic Properties of Individual Species}.
\newblock NASA/TP-2002-211556, 2002. \url{https://ntrs.nasa.gov/citations/20020085330}.

\bibitem{vanOijen2000}
J. A. van Oijen y L. P. H. de Goey.
\newblock Modelling of premixed laminar flames using flamelet-generated manifolds.
\newblock \textit{Combustion Science and Technology}, 161(1):113--137, 2000. DOI: \href{https://doi.org/10.1080/00102200008935814}{10.1080/00102200008935814}.

\bibitem{Bilger1989}
R. W. Bilger.
\newblock The structure of turbulent nonpremixed flames.
\newblock \textit{Symposium (International) on Combustion}, 22(1):475--488, 1989. DOI: \href{https://doi.org/10.1016/S0082-0784(89)80054-2}{10.1016/S0082-0784(89)80054-2}.

\bibitem{Neufeld1972}
P. D. Neufeld, A. R. Janzen y R. A. Aziz.
\newblock Empirical equations to calculate 16 of the transport collision integrals $\Omega^{(l,s)\ast}$ for the Lennard-Jones (12-6) potential.
\newblock \textit{The Journal of Chemical Physics}, 57:1100--1102, 1972.

\end{thebibliography}

\end{document}
